% !TeX root = ../main.tex

\section{Related Work}

We reviewed and discussed the related work with respect to privacy leakage, empirical studies on malware, and malware detection on OSS packages.

% Malicious packages designed to infringe on user privacy pose a significant threat to cybersecurity. These packages are often crafted to secretly extract sensitive information, such as personal identification data, login credentials, or browsing habits, without the user's consent. 
% Attackers may embed keyloggers or data exfiltration scripts within these packages, allowing them to monitor and collect private information over time. Such threats are particularly concerning in unsecured or public networks, where data transmission is more vulnerable to interception and exploitation. 

%introduce various attacks that breach privacy
% Currently, there are   %\todo{research work on extracting PII, password, etc.}
% , focusing on forums that trade credential leaks, phishing kits and off-the-shelf keyloggers that harvest passwords. 
% And they showed how blocking login attempts that fail to match a user's historical login behavior or device profile help mitigate the risk of the three types of attacks mentioned above. 
% Another attack to breach privacy is ransomware. It is a type of malicious software where attackers encrypt the victim's files or lock system access, and then extort ransom in exchange for decryption keys or restoring access. 
% Split learning (SL) is a framework to protect a client's training data. 
%Response methods for data breaches
% This method is to launch a false information attack by returning a large amount of bait information to the attacker when there is suspicion of unauthorized access and then using questioning questions for verification.
% provides an enhanced layer of protection by 
% For cyberattacks aimed at breaching privacy, multiple studies have provided methods for detecting malicious software and user response strategies. 


\subsection{Privacy Leakage}

Many studies investigated cyberattacks targeting privacy leakage. Thomas et al.~\cite{thomas2017data} presented the first longitudinal measurement study of how miscreants obtain stolen credentials and subsequently bypass risk-based authentication schemes to hijack a victim's account. They develop an automated framework that monitors credential thefts in blackmarket. Through this framework, they identified over 14.3 million victims in total between 2016 and 2017, demonstrating the vast underground economy surrounding credential theft. 
Ransomware is another threat to privacy leakage. Gabriela et al.~\cite{almeida2023analyzing} look into the evolving dynamics of ransomware. They analyze network traffic patterns of ransomware activities and found that ransomware is shifting towards complex data leakage strategies. Ozturk et al.~\cite{ozturk2024dynamic} conducted a comprehensive analysis of ransomware variants that specialize in data theft, emphasizing the dynamic behavior patterns and techniques adopted by these malicious entities to compromise data privacy. Gao et al.~\cite{gao2023pcat} propose a novel pseudo-client attack mechanism on various split learning models, revealing serious privacy risks of split learning. 
As for defensive measures, Hirano et al.~\cite{hirano2009t} proposed T-PIM mechanism that provides a secure password input method to users, which features a hypervisor to isolate a trusted domain and provides users with a secure password input method. Chowdhury et al.~\cite{chowdhury2017protecting} proposes an efficient scheme for protecting sensitive data from malicious threats using data mining and machine learning techniques. 
Stofol et al.~\cite{stolfo2012fog} propose a different approach for securing data in the cloud using offensive decoy technology and protect against the misuse of the user's real data. Olabim et al.~\cite{olabim2024differential} proposed a novel framework integrating differential privacy to ensure that exfiltrated data remains unusable to attackers through the introduction of statistical noise.


\subsection{Empirical Studies on Malware.} 
% At present, empirical studies on malware covers multiple aspects, empirical studies on malware aiming to understand its behavior, propagation mechanisms, defense strategies, and impact on society.
% Open source software supply chain attack refers to attackers who tamper with or contaminate the supply chain links of open source software, implant malicious code into legitimate software, and thereby affect downstream users. 
% This type of attack has significantly increased in recent years, and due to its strong concealment and wide impact, it has become an important threat to network security.

Malware empirical studies contribute to the understanding of its behavior, propagation mechanisms, defense strategies, and impact on society. The first systematic investigation into malicious packages was conducted by Ohm et al.~\cite{ohm2020backstabber}. They collected 174 malicious software packages from multiple platforms and manually studied their injection methods, triggering methods, and targets. Guo et al.~\cite{guo2023empirical} investigated the characteristics and current state of the malicious code lifecycle in the PyPI ecosystem, reflecting the characteristics of malicious code at different stages.
Zahan et al.~\cite{zahan2022weak} conducted an empirical study on npm package metadata, they identified six signals of security weakness in NPM supply chain and presented a framework that helps prioritization and quantification of weak link signals in the npm supply chain. 

% With the rapid evolution of malware technology, detection techniques are also constantly innovating. 

\subsection{Malware Detection on OSS Packages.} 
Malware Detection is one of the core areas of cybersecurity, aimed at identifying, classifying, and preventing attacks by malicious software (such as ransomware, spyware, etc.) through technological means.
Garrett et al.~\cite{garrett2019detecting} selected multiple features and used clustering techniques to construct a benign behavior model for detecting malicious NPM packages.
Ohm et al.~\cite{ohm2022feasibility} extended their work on this basis. They added features related to software package metadata and obfuscation and evaluated the feasibility of different supervised ML models in detecting malicious packages.
Sejfia et al.~\cite{sejfia2022practical} present Amalfi, a ML based approach for automatically detecting potentially malicious packages comprised of three complementary techniques. In this work, they used the feature set extended by Garrett et al.~\cite{garrett2019detecting} to train the classifier. There are also many tools for malware detection in the PyPI ecosystem. The detection method proposed by Vu et al.~\cite{vu2020typosquatting} focuses on typosquatting and combosquatting attacks. They perform detection by checking whether the name of the package is similar to the package in the Python standard library or a known source code repository. Liang et al.~\cite{liang2023needle} proposed using clustering techniques to group PyPI installation scripts for detection. And Sun et al.~\cite{sun2024integrating} combined deep code behavior features with metadata features for detection. Both of them model malicious behavior as discrete features.

% zhang et al. , sun et al.

% zhang2023malicious sun2024integrating
